\section{Video Classification}
\label{sec:app_video_classification}

We sample 128 and 32 frames with a stride of 2 frames from Kinetics 400 videos~\citep{kay2017kinetics} and Moments in Time~\citep{monfort2019moments} videos, respectively. For both ViT-22B and ViT-e we rely in the frozen, pre-trained models and use the pre-logit feature representation to extract a single embedding per frame, resulting in a token sequences of length 128 and 32, respectively, which are then processed by a shallow transformer model equipped with a class-token classifier.

This is in contrast to CoCa~\citep{yu2022coca}, which uses one token per image patch for their video classification experiments and a resolution of 576px (compared 224px in our experiments), resulting in much longer token sequences. We explored using one token per image patch (i.e. unpooled features) in preliminary experiments, but found that this leads to inferior performance. One potential reason for this could be that CoCa applies a contrastive loss to a pooled feature representation, and additionally feeds the unpooled token sequences to a generative decoder, which might lead to a different structure in the unpooled representation than the supervised classification loss used to pretrain ViT-22B and ViT-e.

To facilitate experimentation, we pre-compute frozen features for the two ViT variants we consider, using the same augmentations as~\citep{arnab2021vivit}. To improve the robustness of our model and prevent overfitting we feed the entire training set ten times, with different data augmentations for every pass. We train for 30 epochs on these precomputed features with a batch size of 256 using SGD with momentum and with a cosine schedule including a linear warmup of 2.5 epochs. We sweep the following hyperparameters and corresponding value ranges to train our video model: $\{1, 2\}$ transformer layers of width $\{1024, 2048, 4096\}$, using a learning rate in $\{10^{-1}, 10^{-2}\}$ and a weight decay in $\{10^{-3}, 10^{-2}\}$.
